\documentclass[twoside]{report}
\usepackage[czech]{babel}
\usepackage{fancyhdr}
\usepackage{graphicx}
\usepackage{fontspec}
\usepackage{titling}
\usepackage[sfdefault]{roboto} % Roboto - jde na něj koukat, a funguje ve Windows i Linuxu
\usepackage[setpagesize=false, 
            unicode=false, 
            colorlinks=true, 
            urlcolor=blue, 
            linkcolor=black]{hyperref} 
\usepackage[a4paper, 
            left=2cm,
            right=2cm,
            top=2.5cm,
            bottom=2.5cm,
            footskip=1.25cm]{geometry}
\usepackage{longtable}
\usepackage{booktabs}
\usepackage{array}
\usepackage{multirow}
\usepackage[table]{xcolor}
\usepackage{wrapfig}
\usepackage{float}
\usepackage{colortbl}
\usepackage{pdflscape}
\usepackage{tabu}
\usepackage{threeparttable}




\setlength{\parindent}{0pt}
\setlength{\parskip}{1em plus 2pt minus 1pt}
\setlength{\emergencystretch}{3em} % prevent overfull lines

\newcommand{\Nazev}{$if(title)$ {$title$} $else$ {Bezejmenný report} $endif$} % Název reportu (nebo hláška, že není)

\newcommand{\Paticka}{$if(project)$ {$project$} $endif$ $if(version)$ { $version$} $endif$ $if(date)$ { $date$} $endif$} % patička - když projekt tak projekt, když verze tak verze, když datum tak datum; mezi tím mezera

\renewcommand{\baselinestretch}{1.05} % menší, ale platné mezery mezi řádky 

\providecommand{\tightlist}{\setlength{\parskip}{0pt}} % definovat \tightlist, aby pandoc nepadal...

\pagestyle{fancy}

\makeatletter % bez číslování nadpisů (sekcí)
\renewcommand{\@seccntformat}[1]{}
\makeatother

\addto\captionsczech{\renewcommand{\figurename}{Obr.}} % Obrázek je pro děti...
\raggedbottom % aby obrázek nerozhodil předchozí stránku...

\begin{document}

\fancyhf{} % sets both header and footer to nothing
\renewcommand{\headrulewidth}{0pt} % pryč se zbytečnou čárou!

\fancyfoot[RE,RO]{\textit\thepage} % vnější patička - číslo stránky
\fancyfoot[CE,CO]{} % prostřední patička - prázdná
\fancyfoot[LE,LO]{\includegraphics[scale=0.02]{/home/jindra/RPTG/iota.pdf}\textit{ JLA} \textit\Paticka} % vnitřní patička - verze a datum v závorce

\begin{center} % titulek s podtržítkem
\fontseries{bx}{\LARGE\Nazev}
\noindent\hfil\rule{.8\textwidth}{.25pt}\hfil % podtržení pod titulkem, ne přes celou stranu a ne moc tlusté
\end{center}

$body$ % vlastní text z Markdownu

\end{document} 
